\scanpage{16}
- \mspage{13} -
\uline{sans fibre additive}, \uline{munies de formes différentielles invariantes} $\pi$ \uline{et} $\pi'$, \uline{et de même invariant modulaire}. \uline{On suppose que} 2, $c_4$ \uline{et} $c_6$ \uline{sont inversibles}, \uline{i.e.} \uline{que} $2^{-1}$, $1/j$ et $1/(j-1728)$ \uline{sont des sections de} $S$. \uline{Alors}, \uline{le schéma} $\uline{\mathrm{Isom}}(E,F)$ \uline{s'identifie au schéma défini par l'équation}
\[
\lambda^2=c_4c'_6/c_6c'_4.
\]
En vertu de (3.6), se donner $(E,\pi)$ revient à se donner $a_2$, $a_4$ et $a_6$ vérifiant $a_2a_4-9a_6=0$, et un isomorphisme de $E$ sur $E'$ est décrit par $\lambda$ tel que $(E\ \pi)$ soit isomorphe à $(E',\lambda\pi')$, i.e. par $\lambda$ vérifiant
\[
(6.4.1)\qquad a'_2=\lambda^2a_2\qquad a'_4=\lambda^4a_4\qquad a'_6=\lambda^6a_6.
\]
Par hypothèse, $j=j'$, de sorte que $(c_4^3,c_6^2,\Delta)$ est proportionnel à $({c'_4}^3,{c'_6}^2,\Delta')$, et que $\mu=c_4c'_6/{}_{c_6c'_4}$ vérifie
\[
(6.4.2)\qquad c'_4=\mu^2c_4\qquad c'_6=\mu^3c_6\qquad\Delta'=\mu^6\Delta.
\]
En vertu de (3.4), on a alors
\[
a'_2=\mu a_2\qquad a'_4=\mu^2a_4\qquad a'_6=\mu^3a_6.
\]
\margin{\figure{\begin{tikzpicture}\draw (0,0)--(.35,-.015);\end{tikzpicture}}}
De plus, d'après (3.4) encore, localement, soit $a_2$, soit $a_4$ et $a_6$ sont inversibles, de sorte que (6.4.1) équivaut à $\lambda^2=\mu$.

Pour prouver (III) en général, on peut se ramener à supposer $S$ noethérien, et supposer soit 2 inversible (cas traité en (6.4)), soit
\margin{15}
